\documentclass[12pt,twoside,letterpaper]{article}
% Rapport PFE — Mes histoires traduites (StoryTranslatorWeb)
% Overleaf : garder includes.tex, uqar.png, kpfonts.sty, etc. du ZIP du cours.
% Remplacer ce fichier + report_v1.0.bib + titlepage.tex, puis ajouter figures/

\newcommand{\reporttitle}{{\huge \textsc{Rapport} \\ \vspace{0.5cm} {\Large Projet de fin d'études}}}

\newcommand{\reportauthorOne}{Aissatou Seck}
\newcommand{\reportauthorTwo}{}
\newcommand{\reportauthorThree}{}
\newcommand{\reportProfessor}{Mohamed Tarik Moutacalli}
\newcommand{\reporttype}{Coursework}
\bibliographystyle{plain}

% Nom commercial de l'application (affiché dans l'interface)
\newcommand{\appname}{Mes histoires traduites}
\newcommand{\appcode}{StoryTranslatorWeb}

\input{includes}

% Encadré si une capture n'est pas encore déposée dans figures/
\newcommand{\captureecran}[3]{%
\begin{figure}[H]
\centering
\IfFileExists{#1}{%
  \includegraphics[width=0.92\textwidth]{#1}%
}{%
  \fbox{\parbox{0.9\textwidth}{\centering\vspace{2.2cm}
  \textbf{Capture d'écran à ajouter}\\[0.4cm]
  Fichier~: \texttt{#1}\\[0.2cm]
  #2\vspace{2.2cm}}}%
}%
\caption{#3}
\label{fig:#2}
\end{figure}
}

%%%%%%%%%%%%%%%%%%%%%%%%%%%%

\begin{document}
% Page de titre — adapter au besoin (logo uqar.png déjà dans le ZIP Overleaf)
\begin{titlepage}

\newcommand{\HRule}{\rule{\linewidth}{0.5mm}}

\begin{center}

\includegraphics[width=7cm]{uqar}\\[1.5cm]
\textbf{\textsc{\Large Projet de fin d'études en informatique \\ \vspace{0.5cm} INF39615-7T}}\\[1.0cm]
\textsc{\Large Université du Québec à Rimouski}\\[0.5cm]
\textsc{\large Département de mathématiques, d'informatique et de génie}\\[0.95cm]

\HRule \\
{ \huge \bfseries \reporttitle}\\
\HRule \\[1.5cm]
\end{center}

\begin{flushleft} \large
\textit{\textsc{\textbf{Équipe :}}}\\
\reportauthorOne \\
\vspace{0.5cm}
\textit{\textsc{\textbf{Professeur :}}}\\
\reportProfessor
\end{flushleft}

\vfill
\makeatletter
Date: \@date
\makeatother

\end{titlepage}


\tableofcontents
\newpage

%%%%%%%%%%%%%%%%%%%%%%%%%%%%

\section{Introduction}
\label{sec:introduction}

\appname{} (dépôt \appcode{}~\cite{ref_github_projet}) est une application web destinée aux enfants qui souhaitent apprendre une langue étrangère grâce à la lecture d'histoires illustrées.
L'idée pédagogique est simple~: l'enfant lit une histoire scène par scène, peut l'écouter à voix haute, change éventuellement la langue d'affichage du récit, puis vérifie sa compréhension à travers un quiz guidé par une intelligence artificielle.

Une première version avait été réalisée dans le cadre du cours INF34615-7T.
Elle permettait déjà la lecture multilingue, la synthèse vocale du navigateur et une forme de sauvegarde de progression.
Le présent rapport porte sur la \textbf{continuation} du projet en \textbf{INF39615-7T}, sous la supervision du professeur Mohamed Tarik Moutacalli.

Les objectifs principaux de cette continuation étaient les suivants~:
\begin{itemize}
    \item ajouter des \textbf{comptes utilisateurs} sécurisés (inscription, connexion, rôles)~;
    \item assurer la \textbf{persistance} des données en base PostgreSQL (profil, favoris, progression, historique)~;
    \item distinguer un \textbf{espace enfant} et un \textbf{espace administrateur}~;
    \item adapter l'\textbf{interface} à la langue maternelle de l'utilisateur~;
    \item proposer un \textbf{quiz pédagogique} en fin d'histoire, avec correction linguistique~;
    \item permettre un \textbf{mode oral} fondé sur la transcription Whisper~\cite{ref_whisper_paper,ref_whisper_api}.
\end{itemize}

Les sections suivantes décrivent le fonctionnement du site (avec captures d'écran), les technologies utilisées, puis les éléments techniques utiles à la maintenance et à l'évolution du projet.

\section{Présentation du projet}
\label{sec:presentation}

Cette section décrit le parcours utilisateur tel qu'il apparaît dans l'application.
Chaque sous-section correspond à une fonctionnalité majeure, illustrée par une capture lorsqu'elle est disponible dans le dossier \texttt{figures/}.

\subsection{Page d'accueil publique}
\label{sec:accueil}

Avant toute connexion, le visiteur arrive sur une page d'accueil moderne qui présente la marque \appname{}, un message clair sur l'apprentissage des langues par les histoires, et des actions principales~: découvrir le site, créer un compte ou se connecter.
Le bouton \emph{Découvrir} ouvre une section écrite qui explique les atouts de la plateforme (histoires interactives, langues, espace personnel, quiz avec l'IA), plutôt que d'afficher directement le catalogue d'histoires.
Un pied de page regroupe la navigation, les informations légales, les coordonnées et un formulaire d'inscription à la lettre d'information.

\captureecran{figures/01_accueil.png}{accueil}{Page d'accueil publique de \appname{}.}

\subsection{Authentification et comptes utilisateurs}
\label{sec:auth}

Avant d'accéder à la bibliothèque, l'utilisateur doit créer un compte ou se connecter.
À l'inscription, l'enfant (ou l'adulte qui l'accompagne) renseigne notamment~: prénom, nom, courriel, langue maternelle, langue cible, niveau, tranche d'âge et mot de passe.

Les mots de passe sont hachés avec \textbf{bcrypt} et ne sont jamais stockés en clair.
Les sessions s'appuient sur des jetons \textbf{JWT}~\cite{ref_jwt} transmis au client après une authentification réussie.
Deux rôles existent dans le système~:
\begin{itemize}
    \item \texttt{user}~: espace enfant (bibliothèque filtrée par âge, profil pédagogique, lecture, quiz)~;
    \item \texttt{admin}~: espace administrateur (gestion de la bibliothèque, import d'histoires, liste des utilisateurs).
\end{itemize}

\captureecran{figures/02_inscription.png}{inscription}{Formulaire d'inscription (création de compte).}

\captureecran{figures/03_connexion.png}{connexion}{Écran de connexion.}

\subsection{Page d'accueil après connexion~: la bibliothèque}
\label{sec:biblio}

Après authentification, l'utilisateur \texttt{user} arrive dans son espace personnel (espace enfant).
La bibliothèque affiche les histoires disponibles sous forme de cartes (image de couverture, titre, indication de progression).

Depuis cette page, l'utilisateur peut~:
\begin{itemize}
    \item cliquer sur une carte pour lancer la lecture~;
    \item consulter ses \textbf{favoris} et ses histoires \textbf{récemment lues}~;
    \item ouvrir son profil pour modifier ses informations~;
    \item lancer une \textbf{histoire au hasard} parmi celles autorisées pour son âge.
\end{itemize}

Les histoires affichées dépendent de la tranche d'âge déclarée dans le profil~: 3 à 5~ans (\texttt{petits}), 6 à 8~ans (\texttt{moyens}), 9 à 12~ans (\texttt{grands}).
Cette règle pédagogique évite de proposer un contenu trop difficile ou trop simple.

\captureecran{figures/04_bibliotheque.png}{bibliotheque}{Bibliothèque de l'espace enfant.}

\subsection{Profil utilisateur et préférences}
\label{sec:profil}

Dans \emph{Mon profil}, l'utilisateur peut modifier ses informations et sa photo.
Deux préférences linguistiques sont particulièrement importantes~:
\begin{itemize}
    \item la \textbf{langue maternelle} détermine la langue des textes de l'interface (menus, boutons, messages)~;
    \item la \textbf{langue cible} est celle que l'enfant apprend~: c'est dans cette langue que le quiz pose ses questions et attend les réponses.
\end{itemize}

Le changement de langue d'interface est appliqué après l'enregistrement du profil (et non au simple changement du menu déroulant), afin d'éviter des rechargements intempestifs et de garantir une expérience stable.

\captureecran{figures/05_profil.png}{profil}{Écran de profil et préférences linguistiques.}

\subsection{Importation d'une histoire}
\label{sec:import}

L'importation d'histoires est réservée à l'administrateur.
Celui-ci clique sur \emph{Importer} et sélectionne un fichier \texttt{.zip} contenant~:
\begin{itemize}
    \item un fichier \texttt{story.json} (titre, scènes, textes multilingues, chemins d'images, tranche d'âge)~;
    \item les images des scènes (et éventuellement des avatars de personnages).
\end{itemize}

Le serveur décompresse le fichier, vérifie la présence et la validité de \texttt{story.json}, puis ajoute l'histoire à la bibliothèque.
Des contrôles de sécurité limitent la taille du ZIP, le nombre de fichiers et le volume décompressé (protection contre les archives malveillantes).
Si le fichier est invalide, il est rejeté avec un message d'erreur clair.

\captureecran{figures/06_import_admin.png}{import}{Action d'importation d'une histoire (espace administrateur).}

\subsection{Espace administrateur}
\label{sec:admin}

L'administrateur dispose d'un tableau de bord distinct, avec un accès bien visible à l'import ZIP et une vue d'ensemble des catégories d'âge.
Son profil est volontairement simplifié~: prénom, nom, courriel et photo uniquement (pas de paramètres pédagogiques enfant).
Une section permet également de consulter la liste des utilisateurs inscrits, utile pour le suivi et la démonstration du projet.

\captureecran{figures/07_admin.png}{admin}{Tableau de bord administrateur.}

\subsection{Le lecteur d'histoire}
\label{sec:lecteur}

Une fois une histoire sélectionnée, l'interface bascule en mode lecteur.
L'utilisateur voit~:
\begin{itemize}
    \item l'image de la scène (plein écran, adaptée à la hauteur de la fenêtre)~;
    \item l'avatar du personnage~;
    \item la bulle de texte dans la langue sélectionnée~;
    \item une barre et un compteur de progression (scène courante / total)~;
    \item les boutons de navigation (précédent / suivant)~;
    \item sur la dernière scène, un bouton pour discuter avec l'IA.
\end{itemize}

Le lecteur a été conçu pour que \textbf{tous les contrôles restent visibles sans faire défiler la page}, quel que soit l'utilisateur connecté et quelle que soit l'histoire ouverte.
L'image occupe le fond de l'écran~; les boutons sont superposés de façon stable.
Si le texte d'une scène est très long, seul le contenu de la bulle peut défiler, pas la page entière.

\captureecran{figures/08_lecteur.png}{lecteur}{Lecteur d'histoire (scène, texte et contrôles).}

\subsubsection{Changement de langue}
Un contrôle permet de faire tourner la langue d'affichage du texte de l'histoire parmi les sept langues supportées (français, anglais, arabe, espagnol, allemand, italien, portugais).
Lorsque l'arabe est sélectionné, le texte s'affiche de droite à gauche (\texttt{rtl}).

\subsubsection{Contrôles audio}
Le site utilise la synthèse vocale du navigateur (Web Speech API~\cite{ref_speech}) pour lire le texte à voix haute~:
\begin{itemize}
    \item lecture automatique à chaque changement de scène (activable / désactivable)~;
    \item lecture manuelle (lecture / pause / reprise)~;
    \item indication de la langue de lecture au navigateur (ex.\ \texttt{fr-FR}, \texttt{en-US}, \texttt{ar-SA}).
\end{itemize}

\subsubsection{Mode immersif (masquer le texte)}
Un bouton permet de masquer la bulle de texte afin de favoriser la compréhension orale~: l'enfant écoute d'abord, puis peut réafficher le texte pour vérifier.

\subsubsection{Sauvegarde de progression}
La progression est sauvegardée en base de données pour l'utilisateur connecté.
À la réouverture d'une histoire déjà commencée, une boîte de dialogue propose de reprendre à la scène enregistrée ou de recommencer depuis le début.

\subsection{Quiz de fin d'histoire avec l'IA}
\label{sec:quiz}

À la fin d'une histoire, l'enfant peut discuter avec une IA (Gemini~\cite{ref_gemini}) qui~:
\begin{itemize}
    \item pose des questions uniquement sur le contenu de l'histoire~;
    \item demande éventuellement un court résumé~;
    \item corrige les erreurs de langue de manière pédagogique~;
    \item communique exclusivement dans la \textbf{langue cible}~;
    \item refuse les chiffres dans les réponses (l'enfant doit écrire ou dire les nombres en toutes lettres).
\end{itemize}

Cette activité transforme la lecture en un véritable exercice de compréhension et d'expression, plutôt qu'en un simple parcours passif d'images.

\captureecran{figures/09_quiz.png}{quiz}{Quiz conversationnel en fin d'histoire.}

\subsection{Mode oral et transcription Whisper}
\label{sec:whisper}

Le quiz propose un mode \emph{Écrire} et un mode \emph{Parler}.
En mode Parler~:
\begin{enumerate}
    \item l'IA lit sa question à voix haute~;
    \item l'enfant enregistre sa réponse au micro (MediaRecorder~\cite{ref_mediarecorder})~;
    \item l'audio est envoyé au serveur~;
    \item Whisper (OpenAI)~\cite{ref_whisper_api,ref_whisper_paper} transcrit la voix de façon littérale~;
    \item le texte est transmis à l'IA pour correction et suite du quiz.
\end{enumerate}

Cette fidélité de transcription est essentielle~: l'IA doit voir les vraies erreurs pour pouvoir les corriger.
Le mode écrit reste disponible si le micro n'est pas autorisé ou si l'enfant préfère taper sa réponse.

\captureecran{figures/10_quiz_oral.png}{quizoral}{Quiz en mode oral (micro et transcription).}

\subsection{Interface multilingue (langue maternelle)}
\label{sec:i18n}

Les menus, boutons et messages suivent la langue maternelle du profil, y compris après déconnexion puis reconnexion.
Le système de traduction côté client repose sur un fichier de messages et une fonction \texttt{t()} (ou \texttt{translate()}) qui sélectionne la chaîne selon la locale.
Dans les listes de langues, chaque langue apparaît dans sa forme native (ex.\ English, Español, العربية, Deutsch) pour être facilement identifiable, quelle que soit la langue d'interface.

\newpage

\section{Technologies utilisées}
\label{sec:technologies}

\subsection{Front-end (interface client)}
\begin{itemize}
    \item \textbf{React.js}~\cite{ref_react}~: construction de l'interface et organisation en composants.
    \item \textbf{Vite}~\cite{ref_vite}~: outil de développement et de compilation du client.
    \item \textbf{Tailwind CSS}~\cite{ref_tailwind}~: mise en forme responsive.
    \item \textbf{Axios}~: communication HTTP avec l'API.
    \item \textbf{Lucide-react}~: icônes de l'interface.
    \item \textbf{Web Speech API}~\cite{ref_speech}~: synthèse vocale.
    \item \textbf{MediaRecorder}~\cite{ref_mediarecorder}~: enregistrement audio pour Whisper.
\end{itemize}

\subsection{Back-end (serveur)}
\begin{itemize}
    \item \textbf{Node.js}~\cite{ref_nodejs} et \textbf{Express}~\cite{ref_express}~: API REST.
    \item \textbf{Multer} / \textbf{Unzipper}~: réception et extraction des archives ZIP.
    \item \textbf{CORS}~: autorisation contrôlée des appels client~$\leftrightarrow$~serveur.
    \item \textbf{JWT}~\cite{ref_jwt} et \textbf{bcrypt}~: authentification et mots de passe.
\end{itemize}

\subsection{Base de données et authentification}
\begin{itemize}
    \item \textbf{PostgreSQL}~\cite{ref_postgres}~: persistance des comptes, favoris, progression, historique et préférences.
    \item \textbf{Prisma}~\cite{ref_prisma}~: schéma et migrations (la logique runtime du projet utilise principalement le client \texttt{pg} pour les requêtes SQL).
\end{itemize}

Les histoires elles-mêmes sont également gérées de manière structurée (métadonnées et scènes), afin de faciliter le filtrage par âge et la maintenance du catalogue.

\subsection{Intelligence artificielle et reconnaissance vocale}
\begin{itemize}
    \item \textbf{Gemini}~\cite{ref_gemini}~: quiz conversationnel et correction linguistique.
    \item \textbf{Whisper (OpenAI)}~\cite{ref_whisper_api}~: transcription orale littérale.
\end{itemize}

\subsection{Qualité et tests}
\begin{itemize}
    \item \textbf{Vitest}~\cite{ref_vitest}~: tests unitaires du projet (client et serveur).
    \item \textbf{GitHub Actions}~\cite{ref_gha}~: intégration continue (CI). À chaque \emph{push} ou \emph{pull request} pertinente, les tests sont exécutés automatiquement pour détecter les erreurs rapidement. Il n'y a pas, à ce stade, de déploiement automatique (pas de CD).
\end{itemize}

\newpage

\section{Description technique pour la maintenance}
\label{sec:maintenance}

Cette section est destinée à tout développeur souhaitant comprendre, modifier ou faire évoluer le projet.

\subsection{Structure des fichiers}

\begin{verbatim}
StoryTranslatorWeb/
  client/
    src/
      App.jsx
      components/
      i18n/
      utils/
      constants/
      hooks/
  server/
    index.js
    routes/
    services/
    prisma/
    uploads/
  tests/
  .github/workflows/
\end{verbatim}

\subsection{Le fichier client~: \texttt{App.jsx}}
\texttt{App.jsx} reste le point central de navigation entre les vues~: accueil public, authentification, bibliothèque (enfant ou admin), lecteur, quiz.
Le lecteur et la logique de synthèse vocale ont été extraits dans des modules dédiés (\texttt{StoryPlayer.jsx}, \texttt{useStorySpeech.js}) afin de garder le fichier principal plus lisible.

\subsection{Le fichier serveur~: \texttt{index.js}}
Le serveur Express branche plusieurs routes protégées, notamment~:
\begin{itemize}
    \item \texttt{/api/auth}~: inscription, connexion, profil, avatar~;
    \item \texttt{/api/library}~: favoris, historique, progression~;
    \item \texttt{/api/ai}~: début et suite du quiz~;
    \item \texttt{/api/speech}~: transcription Whisper~;
    \item \texttt{/api/stories} et upload ZIP~: bibliothèque d'histoires~;
    \item \texttt{/api/admin}~: fonctionnalités réservées à l'administrateur.
\end{itemize}

\subsection{Le format d'une histoire~: \texttt{story.json}}
Le format initial est conservé. Exemple simplifié~:

\begin{verbatim}
{
  "id": "lina-parc",
  "title": "Une journee au parc",
  "thumbnail": "images/scene1.png",
  "ageBand": "petits",
  "scenes": [
    {
      "id": 1,
      "image": "images/scene1.png",
      "character": {
        "name": "Lina",
        "avatar": "images/scene1.png"
      },
      "text": {
        "fr": "Bonjour ! Je m'appelle Lina.",
        "en": "Hello! My name is Lina.",
        "ar": "..."
      }
    }
  ]
}
\end{verbatim}

\subsection{Comment lancer le projet}
Prérequis~: Node.js et PostgreSQL configurés, fichier \texttt{server/.env} renseigné.
\begin{enumerate}
    \item Terminal~1~: \texttt{cd server}, puis \texttt{npm install}, puis \texttt{node index.js}.
    \item Terminal~2~: \texttt{cd client}, puis \texttt{npm install}, puis \texttt{npm run dev}.
\end{enumerate}
Ouvrir l'URL Vite (par exemple \texttt{http://localhost:5173}).

\subsection{Comment modifier le site}
\begin{itemize}
    \item \textbf{Design}~: classes Tailwind dans les composants React.
    \item \textbf{Nouvelle langue d'interface}~: clés dans \texttt{client/src/i18n/messages.js}.
    \item \textbf{Nouvelle histoire}~: ZIP conforme à \texttt{story.json}, import admin.
    \item \textbf{Prompt du quiz}~: \texttt{server/services/aiChat.js}.
    \item \textbf{Transcription}~: \texttt{server/services/whisperTranscribe.js}.
\end{itemize}

\subsection{Configuration des variables d'environnement}
Le fichier \texttt{server/.env} contient notamment~:
\begin{itemize}
    \item \texttt{DATABASE\_URL}
    \item \texttt{JWT\_SECRET}
    \item \texttt{GEMINI\_API\_KEY}
    \item \texttt{OPENAI\_API\_KEY} (Whisper)
    \item \texttt{WHISPER\_MODEL} (optionnel)
\end{itemize}
Ces secrets ne doivent \textbf{pas} être commités dans Git.

\newpage

\section{Conclusion}
\label{sec:conclusion}

\subsection{Bilan du projet}
L'application \appname{} est passée d'un simple outil de lecture multilingue à une plateforme plus complète.
Elle intègre désormais des comptes sécurisés, une persistance des données, des rôles utilisateur et administrateur, une interface adaptée à la langue maternelle, ainsi qu'un quiz guidé par l'IA avec un mode oral fondé sur Whisper.
L'architecture client/serveur initiale a été conservée et étendue de manière cohérente.

Le dépôt GitHub du projet est accessible ici~: \url{https://github.com/Seck2000/StoryTranslatorWeb}.

\subsection{Ce que ce projet m'a appris}
Ce projet m'a permis de consolider~:
\begin{itemize}
    \item React et la gestion d'état~;
    \item l'architecture client/serveur et les API REST protégées~;
    \item PostgreSQL et la modélisation des données utilisateur~;
    \item l'intégration de Gemini et Whisper~;
    \item les tests unitaires et l'intégration continue (CI) avec GitHub Actions~;
    \item l'importance de l'expérience utilisateur (lecteur sans défilement, i18n, rôles).
\end{itemize}

\subsection{Améliorations futures}
\begin{itemize}
    \item déploiement en ligne (frontend, backend, base hébergée, HTTPS)~;
    \item renforcement des tests d'intégration~;
    \item éditeur d'histoires en ligne~;
    \item amélioration continue de l'expérience mobile~;
    \item éventuel quota d'usage des API IA pour maîtriser les coûts en production.
\end{itemize}

\newpage
\bibliography{report_v1.0}

\end{document}
